%% LyX 2.4.4 created this file.  For more info, see https://www.lyx.org/.
%% Do not edit unless you really know what you are doing.
\documentclass[english]{article}
\usepackage[T1]{fontenc}
\usepackage[latin9]{inputenc}
\usepackage{units}
\usepackage{enumitem}
\usepackage{amsmath}
\usepackage{amsthm}
\usepackage{cancel}
\usepackage{geometry}
\geometry{verbose,tmargin=1in,bmargin=1in,lmargin=1in,rmargin=1in}

\makeatletter

%%%%%%%%%%%%%%%%%%%%%%%%%%%%%% LyX specific LaTeX commands.
%% Because html converters don't know tabularnewline
\providecommand{\tabularnewline}{\\}

%%%%%%%%%%%%%%%%%%%%%%%%%%%%%% Textclass specific LaTeX commands.
\numberwithin{equation}{section}
\numberwithin{figure}{section}
\newlength{\lyxlabelwidth}      % auxiliary length 
\theoremstyle{definition}
\newtheorem*{example*}{\protect\examplename}

%%%%%%%%%%%%%%%%%%%%%%%%%%%%%% User specified LaTeX commands.
\usepackage{fancyhdr}
\usepackage{lastpage}
\pagestyle{fancy}
\usepackage{enumitem}
\usepackage{ifthen}
%sets math fonts to be more readable
\everymath=\expandafter{\the\everymath\displaystyle}
%\DeclareMathSizes{10}{10}{10}{10}
\usepackage{multicol}

\usepackage{tikz}
\usetikzlibrary{plotmarks}
\usepackage{pgfplots}
\pgfplotsset{compat=newest}

\usepackage{calc}
\usepackage[nomessages]{fp}

\usepackage{advdate}    % Advancing/saving dates
\usepackage{datetime}   % Dates formatting
\usepackage{datenumber} % Counters for dates
% Added by lyx2lyx
\setlength{\parskip}{\smallskipamount}
\setlength{\parindent}{0pt}

\makeatother

\usepackage{babel}
\providecommand{\examplename}{Example}

\begin{document}
\renewcommand{\author}{Dr.\,James Ashe}

\newcommand{\classNum}{131}

\newcommand{\classSec}{B}

\newcommand{\nameType}{}

\newcommand{\examType}{Exam 2 Review}

\renewcommand{\date}{\formatdate{28}{10}{2013}}

%%First page's header%%%%%%%%%%%%%%%%%%%%%%
\fancypagestyle{firststyle} %%header settings for first pagr
{
	\fancyhead[L]{MTH \classNum{}(\classSec{}) \author{}\\
	\textbf{\examType{}} \date{}}
	\fancyhead[C]{\textbf{\nameType}}
	\setlength{\headsep}{14pt}
}
%%%%%%%%%%%%%%%%%%%%%%%%%%%%%%%%%%%%%%%%%%

%%Next page's header%%%%%%%%%%%%%%%%%%%%%%%
\setlist{nolistsep}
\lhead{\classNum{}(\classSec{})} 
\chead{\textbf{\nameType}} 
\cfoot{\thepage{}~of~\pageref{LastPage}} 
\setlength{\headsep}{5pt} 
%%%%%%%%%%%%%%%%%%%%%%%%%%%%%%%%%%%%%%%%%%%

%%Various settings; see the preamble for other settings%%%%%
%%\setenumerate[0]{leftmargin=12pt,itemindent=0pt,labelwidth=0pt}
\renewcommand{\labelenumi}{\textbf{\arabic{enumi}.}}
\renewcommand{\labelenumii}{\textbf{\alph{enumii}.}}
\thispagestyle{firststyle}
%%%%%%%%%%%%%%%%%%%%%%%%%%%%%%%%%%%%%%%%%%
\newcommand{\xmax}{10}
\newcommand{\xmin}{-10}
\newcommand{\xstep}{0} %must be xmin+n
\newcommand{\ymax}{10}
\newcommand{\ymin}{-10}
\newcommand{\ystep}{0} %must be ymin+n

\newcommand{\Dspace}{\vspace{1cm}}

\textbf{The graph of a quadratic function is given. Determine the
function's equation from a set of options. }\begin{multicols}{2}
\begin{enumerate}
\item Check the direction, $a$.
\item Check the $y$-int.
\item Check the vertex. \columnbreak
\end{enumerate}
\begin{example*}
Determine the function's equation.

\renewcommand{\xmax}{10} \renewcommand{\xstep}{5}
\renewcommand{\ymax}{\xmax} \renewcommand{\ystep}{5}
%%%%%%%
\begin{tikzpicture} 
	\begin{axis}[
		axis x line=middle, axis y line=middle,     		
		xtick={-\xmax,-\xstep,...,\xmax},
		ytick={-\ymax,-\ystep,...,\ymax},
		minor x tick num={4},
		minor y tick num={4},
		grid=both,
		xlabel={$$},     
		ylabel={$$},
		xmin=-\xmax.25,xmax=\xmax.25,
		ymin=-\ymax.25,ymax=\ymax.25,     		width=6cm,     		height=6cm] 		\addplot[red,very thick,mark=noner,smooth,domain=\xmin:\xmax]{x^2+2*x+1}; 
	\end{axis}
\end{tikzpicture}

\begin{multicols}{2}

\begin{enumerate}
\item [\textbf{(A)}]\textbf{$x^{2}+2x+1$}
\item [\textbf{(B)}]\textbf{$x^{2}-2x+1$\columnbreak}
\item [\textbf{(C)}]\textbf{$x^{2}-1$}
\item [\textbf{(D)}]$-x^{2}-1$
\end{enumerate}
\end{multicols}
\end{example*}
\end{multicols}\vfill

\textbf{We write quadratic functions in one of two ways; $f(x)=a(x-h)^{2}+k$
or $f(x)=ax^{2}+bx+c$. }\begin{multicols}{2}

$f(x)=a(x-h)^{2}+k$
\begin{itemize}
\item Vertex: $(h,k)$
\item $a<0\Rightarrow$points down\\
$a>0\Rightarrow$points up
\item $x$-int: $a(h)^{2}+k$\columnbreak
\end{itemize}
\textbf{$f(x)=ax^{2}+bx+c$}
\begin{itemize}
\item Vertex: $\left(-\nicefrac{b}{2a},f(-\nicefrac{b}{2a})\right)$
\item $a<0\Rightarrow$points down\\
$a>0\Rightarrow$points up
\item $x$-int: \textbf{$f(0)=c$}\columnbreak
\end{itemize}
\end{multicols}\vfill

\textbf{Find the vertex/max/min of a quadratic function. }\begin{multicols}{2}

$f(x)=a(x-h)^{2}+k$
\begin{itemize}
\item Vertex: $(h,k)$
\item Max is $k$ if it points down
\item Min is $k$ if it points up

\begin{example*}
$2(x+3)^{2}-10$. \\
Vertex: $(-3,-10)$. Min: $-10$\columnbreak
\end{example*}
\end{itemize}
\textbf{$f(x)=ax^{2}+bx+c$}
\begin{itemize}
\item Vertex: $\left(-\nicefrac{b}{2a},f(-\nicefrac{b}{2a})\right)$
\item Max is $f(-\nicefrac{b}{2a})$ if it points down
\item Min is $f(-\nicefrac{b}{2a})$ if it points up

\begin{example*}
$2x^{2}+12x-8$. \\
$-\nicefrac{12}{2(2)}=-3$ and $f(-3)=2(-3)^{2}+12(-3)-8=-10$. \\
Vertex: $(-3,-10)$. Min: $-10$
\end{example*}
\end{itemize}
\end{multicols}\vfill

\textbf{Determine whether the function is polynomial function.}

A polynomial has the form: $f(x)=a_{n}x^{n}+a_{n-1}x^{n-1}+\cdots+a_{1}x+a_{0}$. 
\begin{itemize}
\item The graph of a polynomial is \emph{smooth }and \emph{continuous}.
\item All the coefficients ($a_{n},a_{n-1},\dots,a_{0}$) are REAL numbers.
Real numbers include fractions and irrational numbers.
\item All the exponents are POSITIVE INTEGERS; this also means that there
are no $x$'s in the denominators. 
\item [\textbf{ex.}]$f(x)=\sqrt{5}x^{5}-2x^{2}+\pi$; YES, $\sqrt{5}$
and $\pi$ are real numbers.
\item [\textbf{ex.}]$f(x)=x^{5}-x^{\nicefrac{3}{2}}+2$; NO, $\nicefrac{3}{2}$
is not an integer.\vfill\newpage
\end{itemize}
\textbf{Use the Leading Coefficient Test to determine the end behavior
of the polynomial function.}
\begin{itemize}
\item %
\begin{tabular}{|c|c|c|}
\hline 
 & $a_{n}$ is positive & $a_{n}$ is negative\tabularnewline
\hline 
$n$ is even & $\nwarrow\cdots\nearrow$ & $\swarrow\cdots\searrow$\tabularnewline
\hline 
$n$ is odd & $\swarrow\cdots\nearrow$ & $\nwarrow\cdots\searrow$\tabularnewline
\hline 
\end{tabular}
\item $n$ is the largest exponent, and $a_{n}$ is the coefficient next
to the $x$ with the largest exponent.
\item [\textbf{ex.}]$f(x)=-x^{4}+2x^{3}-40x^{2}+1$. $a=-1$ and $n=4$
so 

\begin{itemize}
\item []$f(x)$ falls to the left and falls to the right.
\end{itemize}
\item [\textbf{ex.}]
\item [\textbf{ex.}]$f(x)=4x^{2}+2x^{7}-5x+1$. $a=2$ and $n=7$ so 

\begin{itemize}
\item []$f(x)$ falls to the left and rises to the right.\vfill
\end{itemize}
\end{itemize}
\textbf{Find the }\textbf{\emph{x}}\textbf{-intercepts of the polynomial
function. State whether the graph crosses the }\textbf{\emph{x}}\textbf{-axis,
or touches the }\textbf{\emph{x}}\textbf{-axis and turns around.}
\begin{itemize}
\item To find when the \emph{x}-intercepts solve $f(x)=0$. 
\item If a zero has an \emph{even degree }then it turns around. 
\item If a zero has an \emph{odd degree }then it crosses.
\item [\textbf{ex.}]$f(x)=x^{3}-3x^{2}$

\begin{itemize}
\item Solve $f(x)=0$ by factoring $f(x)=x^{2}(x-3)=0$. So then the zeros
are $x=0$ and $x=3$.

\begin{itemize}
\item $x=0$ has degree two; even$\Rightarrow$it turns around.
\item $x=3$ has degree one; odd $\Rightarrow$it crosses.
\end{itemize}
\end{itemize}
\item [\textbf{ex.}]$f(x)=-2(x-1)^{2}(x+2)$

\begin{itemize}
\item $f(x)$ has zeros at $x=1$ and $x=-2$. 

\begin{itemize}
\item $x=1$ has degree two; even$\Rightarrow$it turns around at 1.
\item $x=-2$ has degree one; odd $\Rightarrow$it crosses at $-2$.
\end{itemize}
\end{itemize}
\end{itemize}
\vfill\textbf{Graph the above examples using the intercepts and end
behavior. }
\begin{itemize}
\item $f(x)=x^{3}-3x^{2}$

\begin{itemize}
\item []\renewcommand{\xmax}{10} \renewcommand{\xstep}{5}
\renewcommand{\ymax}{\xmax} \renewcommand{\ystep}{5}
%%%%%%%
\begin{tikzpicture} 
	\begin{axis}[
		axis x line=middle, axis y line=middle,     		
		xtick={-\xmax,-\xstep,...,\xmax},
		ytick={-\ymax,-\ystep,...,\ymax},
		minor x tick num={4},
		minor y tick num={4},
		grid=both,
		xlabel={$$},     
		ylabel={$$},
		xmin=-\xmax.25,xmax=\xmax.25,
		ymin=-\ymax.25,ymax=\ymax.25,     		width=6cm,     		height=6cm] 		\addplot[red,very thick,mark=noner,smooth,domain=\xmin:\xmax]{x^3-3*x^2}; 
	\end{axis}
\end{tikzpicture}\vfill\vspace{4cm}
\end{itemize}
\end{itemize}

\end{document}
